% THE OLD DUKE'S WARS -- CORRECTED GENEALOGY
\documentclass[11pt]{book}
\usepackage[utf8]{inputenc}
\usepackage{geometry}
\geometry{a5paper, margin=1.2in}
\usepackage{graphicx}
\usepackage{epigraph}
\usepackage{lettrine}
\usepackage{xcolor}
\usepackage{titlesec}
\usepackage{fontspec}
\setmainfont{Linux Libertine O}
\setsansfont{Linux Biolinum O}
\setmonofont{Inconsolata}

\titleformat{\chapter}{\huge\bfseries\scshape}{\thechapter}{20pt}{}
\titleformat{\section}{\Large\bfseries}{\thesection}{15pt}{}
\titleformat{\subsection}{\normalsize\bfseries}{\thesubsection}{12pt}{}

\begin{document}

\frontmatter
\title{\Huge\scshape The Old Duke's Wars\\\large Vignettes from the Life of Braedon Fenwood}
\author{Set down by the Gray Wanderer, from the deathbed ledger}
\date{Tarlington, winter of 338 AR}
\maketitle
\tableofcontents

\mainmatter

\chapter*{A Note on This Writing}
\lettrine{T}{his} ledger was found pressed between the leaves of a Canray slate, on the table beside Duke Braedon Fenwood's deathbed. The ink is his own blood -- the last move of a game he could not finish. The pages are foxed, some stained by wine, others by tears. A single braid of Ykrul horsehair marks the place where the old man stopped writing.

The Gray Wanderer has added marginal notes in a drier hand. Where the duke's memory falters, we have bracketed the uncertainty. Where the truth is too cruel for ink, we have left a space.

This is not a chronicle of triumph. It is a record of what a house must bury to survive. And it ends where the next ledger begins -- on the river Belworth, with blood in the water and a crown still warm.

\vspace{1em}
\noindent\textit{--- The Archive of the True Masons, Tarlington chapter, winter of 338 AR}

\clearpage

\chapter{Frame: The Solar at Tarlington}

\lettrine{T}{he} fire burns low. Outside, the Belworth churns with winter melt. Duke Braedon Fenwood -- silver-haired, arthritic, his left arm dead from a wound taken eight years ago at the Sack of Lence -- lies propped on pillows. A Canray slate rests on his chest. The pieces are not set. He moves them in his sleep.

You have been summoned to his solar. Not as servants. As witnesses.

\vspace{1em}
\noindent\textit{"Sit,"} he says. \textit{"I will not rise again. But I will speak. The road does not forgive a debt. It only reschedules it. You are the ledger now."}

\vspace{1em}
\noindent He tells you seven stories. They are not in order. They are not kind. They are his confession.

\clearpage

\chapter{The Black Banner Years (Age Twenty--Two)}

\textbf{Narrator: Braedon Fenwood, recalling his youth. The fire pops. He smiles.}

\lettrine{I}{} was a second son. My brother Conlin was the heir -- clever, careful, born with a club foot that our father called a curse and I called an excuse. I had no land, no title, no purpose. So I ran.

Baldvyn ran with me. He was also a second son -- King Harald IV's second boy, not the heir. Harald the Bold was the elder; Baldvyn and I squired together under Oswin Everblood. We sharpened each other's blades and lied about each other's conquests. When I said, \textit{"I"m leaving,"} he said, \textit{"I"ll bring the horses."}

We joined the Black Banners. Not the condotta of the marches -- the real Banners, the ones who fought the Ykrul raiders in the steppe passes. Our captain was a woman called S?ren the One‑Eyed, Linnic, scarred, and patient as frost. She taught us the first rule of war: \textit{"Count exits, not victims. The road remembers."}

\section*{The First Blood}
Our first skirmish was a ford crossing. I was seventeen and terrified. Baldvyn was eighteen and pretending not to be. The Ykrul came out of the mist with wolf‑tail standards. I killed a man -- a boy, really -- with a spear through the throat. He looked at me as if I had stolen something.

I did. I stole his tomorrows.

Baldvyn pulled me up. \textit{"Don"t name him,"} he said. \textit{"Names are anchors. You"ll carry him forever if you speak it."}

I did not speak it. I still see his face.

\section*{The Bond}
Captain S?ren matched us in the Kon"reh board. She played with stones, not pieces -- river‑smooth, each one a different shade of grey. \textit{"Play the player,"} she said, \textit{"not the position."} Baldvyn beat me three times. I beat him twice. We never kept score. Friendship is not a ledger.

Later, in a camp by the Yloka, we swore an oath. Not on a book -- on a braid of horsehair cut from the mane of a fallen Ykrul mare. \textit{"If one of us becomes king, the other will be his hand. If one of us dies, the other will avenge him. If we both die, the road will remember."}

He became king. I became his hand. I also married his sister, Aobrey. The road remembers.

\section*{The Return}
My brother Conlin died of a fever. I inherited the duchy. I came home to Tarlington with a Ykrul braid in my pocket, a Kon"reh board in my saddlebag, and Aobrey on a pillion behind me. My father met us at the gate. \textit{"You smell of the steppe,"} he said. Then he kissed my bride's hand. \textit{"Good. Stinks of honest work."}

\vspace{1em}
\noindent\textit{Thus endeth the first vignette. The fire cracks. Braedon drinks watered wine.}

\clearpage

\chapter{The Thornhill Rebellion (Age Thirty--Six)}

\textbf{Narrator: Duke Braedon Fenwood, now grey at the temples. He touches a scar on his forearm.}

\lettrine{T}{he} rape of Alayse Everblood -- Baldvyn's youngest sister -- was not a secret. Branton Thornhill made it a boast. He sent the king a lock of her hair wrapped in a marriage contract. \textit{"She came willingly,"} he wrote. \textit{"Your sister is a whore, and I am her lord."}

King Harald IV rode to Brigh. He never came back. They killed him in a ditch, his own sword in his belly.

Then Harald the Bold -- the elder brother, the heir -- stormed the same walls. An arrow took his eye eleven days later. He lingered. He did not recover.

They crowned Baldvyn in a field, with a borrowed coronet and a dead brother's blood on his hands.

\section*{The Siege of Brigh}
Baldvyn summoned his banners. I brought eight hundred Fenwood spears -- every man who could hold a pike, every woman who could ride. We encamped in the shadow of Branton's keep. The walls were high. The garrison was hungry. Branton mocked us from the parapet.

\emph{"Fenwood,"} he called down. \emph{"You were a mercenary. How much did they pay you to betray your people?"}

I did not answer. Baldvyn looked at me. \textit{"He wants rage. Give him geometry."}

I spent the night tracing the keep's drainage sluice. By dawn, we had a tunnel. By noon, we were inside.

\section*{The Duel}
Branton met Baldvyn in the great hall. The fire was out. The tapestries had been cut down for bandages. They fought with longswords -- no shields, no quarter. I held the door against Branton's guard.

Baldvyn was faster. Branton was stronger. It ended when Baldvyn caught Branton's blade on his gorget, shattered the steel, and drove his own sword through the lord's chest.

\emph{"For Alayse,"} he said. \emph{"For Harald. For my brother."} He twisted the blade. Branton died with a smile.

\section*{The Child}
A servant woman fled the keep carrying a boy -- Logan Thornhill, Branton's son, perhaps four years old. Baldvyn drew his dagger. I caught his wrist.

\emph{"He is a child."}

\emph{"He is a seed."}

\emph{"Then prune him gently, or the root will spread."}

Baldvyn lowered the blade. He exiled the boy. The woman took him to Vhasia. Years later, that boy would kill Lyonel Vayle's father, and Veralyn would hunt him to a village called Coldmorrow, and she would burn it to ash.

\emph{"Mercy,"} Baldvyn said later, \emph{"is a debt. You just rescheduled it."}

\vspace{1em}
\noindent\textit{Braedon sets down his cup. His hand trembles. The fire has burned low.}

\clearpage

\chapter{The Longblood Defiance (Age Forty--Eight)}

\textbf{Narrator: The duke's voice is hoarse. He touches a sealed letter -- a queen's plea.}

\lettrine{E}{lianor} of Adria was a warrior queen. She had fought beside Baldvyn, borne him children, ruled the marches when the king was away at war. But Baldvyn grew tired. He took lovers -- men and women both, openly, cruelly. Elianor's pride curdled into rebellion.

She raised the banners of Adria, Ethor, and the Bloods of Lence. Her son Valdais -- or perhaps her brother? the lineage was tangled -- rode at her side. Baldvyn called his host. I sat at the council table and watched two people I loved prepare to murder each other.

\section*{The Failed Mediation}
I rode to Elianor's camp under a flag of parley. She met me in a tent that smelled of wine and grief.

\emph{"You were his friend,"} she said. \emph{"You know what he is."}

\emph{"I know what he was. He is still my king."}

\emph{"Then you will die for him."}

\emph{"I will die for the realm."}

She laughed. It was not a kind sound. \emph{"There is no realm. There are only debts and ledgers."}

I rode back. Baldvyn asked me what she said. I told him. He nodded. \emph{"Then let the ledgers speak."}

\section*{The Battle of Charau}
I commanded the left flank. The ground was wet with autumn rain. The Bloods of Lence charged us three times. Three times we threw them back. Valdais -- young, arrogant, his mother's fury in his eyes -- broke through a gap in the line and rode straight for Baldvyn's standard.

I caught him with a mace. His horse went down. He crawled for his sword. I put my boot on his hand.

\emph{"Yield."}

He spat at me.

I could have killed him. I should have killed him. But I saw Elianor's face in his. I pulled him up and marched him, stumbling, to Baldvyn.

\section*{The Pardon}
Baldvyn wanted execution. \emph{"He raised arms against his king."}

\emph{"He is your son -- or your half‑brother. No one is certain. Does it matter?"}

\emph{"He will rebel again."}

\emph{"Then let him owe his life to mercy. A debt unpaid is a leash."}

Baldvyn chewed his rage for a long minute. Then he ordered Valdais confined to his estates. Elianor was exiled to a convent. The war ended.

Valdais never rebelled again. He became a quiet, loyal duke. And when Baldvyn died, Valdais wept.

\vspace{1em}
\noindent\textit{"Mercy," Braedon says, "is not weakness. It is a long loan."}

\clearpage

\chapter{The Sack of Lence (Age Sixty)}

\textbf{Narrator: Braedon's left arm is dead. He cannot lift it. He looks at it as if it belongs to someone else.}

\lettrine{T}{he} Vhasian king demanded homage. We had lands in Vhasia -- old Valvano holdings, granted by the first Everbloods. We had never paid tribute. He sent tax collectors. I sent them back with their papers shredded.

He declared war. I laughed. Then I raised the banners.

\section*{The Invasion}
My son Seniro commanded the vanguard. Wymund the Red -- hot‑headed, boisterous, secretly gentle. He rode with the Everbloods of Aberiel, with Harald and Elyas beside him. Young knights, hungry for glory. Veralyn was fifteen. She dressed as a page, hid in a supply wagon, presented herself at my tent.

\emph{"I will fight."}

\emph{"Your father would kill me."}

\emph{"My father is not here."}

I assigned her to Seniro as a squire. She watched the streets of Lence run with blood.

\section*{The Siege}
We broke the walls on the third day. Street fighting -- house to house, room to room. I lost my horse to a crossbow bolt. My left arm went dead. A Vhasian spearman knelt over me, raised his weapon.

Seniro cut him down. \emph{"Stay down, Father."}

I stayed down. I watched from the mud as Seniro's men burned the palace.

\section*{The Last Thornhill}
Logan Thornhill -- the boy I had spared -- was a man now. He had fought for Vhasia, risen to captain, taken a new name. Veralyn tracked him to a village called Coldmorrow.

I did not see the killing. I saw the aftermath: smoke rising from a hundred thatched roofs, and Veralyn standing in the square with a mace dripping gore.

\emph{"He insulted my father,"} she said. \emph{"He killed my cousin."}

\emph{"He was a child once."}

\emph{"He was a monster grown."}

I did not argue. I took her aside. \emph{"You will carry this weight forever. That is the price of the blade. Do not let it crush you. Let it make you sharp."}

She wept. I held her. We did not speak of it again.

\vspace{1em}
\noindent\textit{Braedon touches his dead arm. "The boy I spared killed the man I should have been."}

\clearpage

\chapter{The Dying King (Age Sixty--Four)}

\textbf{Narrator: Braedon holds a sealed letter. The wax bears the stag and sun of the Everbloods.}

\lettrine{B}{aldvyn} called me to his bedside. He was wasted, yellowed, his breath sour with the canker that had eaten his stomach. The physicians had bled him white. Nothing helped.

\emph{"I am dying,"} he said.

\emph{"You are old."}

\emph{"I am sixty‑four."}

\emph{"So am I."}

He laughed, then coughed blood. \emph{"Remember the Yloka? The ford?"}

\emph{"I remember."}

\emph{"You saved my life."}

\emph{"You saved mine first."}

He pressed a letter into my hand. \emph{"Open it when I am gone. Not before."}

\emph{"What is it?"}

\emph{"My confession."}

\section*{The Tournament}
Baldvyn died a week later. His son Wyllmot -- the Warrior, the Pious, the Brute -- planned a grand tournament to honor the succession. Veralyn asked to compete. Wyllmot refused. She rode anyway.

She unhorsed three knights: Arin Tymber, my son Lerris, and Ouen Longblood. Ouen's horse died under him. He called her a whore. She dismounted and broke his jaw.

Tylan Vayle -- quiet, patient, a wall of muscle -- defeated her in the quarterfinal. He did not gloat. He offered her his hand. \emph{"You fought well, cousin."}

She took it. She smiled for the first time in months.

\section*{The Letter}
I opened Baldvyn's letter after the tourney. It said only this:

\emph{"Watch over Wyllmot. He has my rage without my restraint. Watch over Veralyn. She has my mother's fire, and it will burn her if no one holds the bucket. You were my brother in all but blood. Do not let my blood destroy your name."}

I burned the letter. The oath remains.

\vspace{1em}
\noindent\textit{"Wyllmot," Braedon says, "was a brute. But he was my king. I kept my oath."}

\clearpage

\chapter{The Seven Week Uprising (Age Sixty--Five)}

\textbf{Narrator: Braedon's voice is thin. He does not look at the fire.}

\lettrine{M}{yrmis} worshippers rose in the mountains. They killed lords, burned synods, declared a holy war against the Unity of Light. Wyllmot marched to storm their citadel -- Eryas, a fortress carved into a cliff.

I counseled patience. \emph{"Let them starve. Winter is coming."}

Wyllmot heard nothing. \emph{"They killed my cousin. They defiled a temple. I will have blood."}

He led the assault himself. A rockslide caught the vanguard. Wyllmot died with a spear through his chest, still swinging his axe.

\section*{The Succession}
Baldvyn Brightstride -- Wyllmot's son, young, untested -- sat the throne. Veralyn rode for the mountains. I stayed in Tarlington with a garrison of old men and wounded.

\emph{"You are not a warrior anymore,"} my son Constano said.

\emph{"I am a duke. I am a strategist. I am a man who has buried too many friends."}

\section*{The Conspiracy}
We learned later that Alaysha Renard -- Wyllmot's widow, Veralyn's mother -- had prayed to Myrmis on the eve of her wedding. She had sacrificed her favorite hound for three daughters and two sons. The goddess granted her wish -- then demanded payment.

Alaysha had not paid. The goddess took Wyllmot.

I warned Veralyn: \emph{"Trust your mother not at all. She serves a god that demands sacrifice."}

Veralyn said nothing. She had already known.

\vspace{1em}
\noindent\textit{"Some debts," Braedon says, "are not rescheduled. They are inherited."}

\clearpage

\chapter{The Death of Brightstride (Two Weeks Ago -- Age Sixty--Six)}

\textbf{Narrator: The duke sits up. He has a fever. His eyes are clear.}

\lettrine{B}{aldvyn} Brightstride was young, handsome, beloved. He had the charm of his father and the restraint of his grandmother. He had been married three months -- to my granddaughter, Jaennis, a waif of a girl with chestnut hair and silent eyes.

He complained of a chill. He drank a posset. He died before dawn.

The physicians said fever. I said poison.

\section*{The Suspects}
The Renards had motive. The Longbloods had opportunity. The Vhasian Regency -- still nursing the wounds of Lence -- had means.

Jaennis wept at the funeral. I watched her. She did not weep for a husband. She wept for a future stolen.

\emph{"Did you kill him?"} I asked her later.

\emph{"I am a Fenwood. I do not answer that question."}

\emph{"Answer it."}

\emph{"He was kind. He was gentle. He did not deserve the crown."}

\emph{"That is not an answer."}

\emph{"It is the only one I have."}

\section*{The Succession Crisis}
Veralyn is the named heir. But Valdais -- Elianor's son, my spared captive -- presses his claim. Aleric, Alaysha's nephew, whispers in dark corners. The realm teeters.

I have written Veralyn a letter. \emph{"Take the crown. Hold it. Let no one tell you that you are not worthy. You are a Fenwood. We are made of exile and iron."}

I have not sent it. I will send it with you.

\vspace{1em}
\noindent\textit{Braedon hands a sealed parchment to the senior PC. "Deliver this. Then come back. The river is rising. The next war begins at dawn."}

\clearpage

\chapter{Epilogue: The Road Reschedules}

\lettrine{T}{he} fire is ash. Braedon Fenwood lies still. His chest rises once, twice -- then stops.

A servant draws the sheet over his face. The Canray slate falls from his hand. The pieces scatter.

One -- a Blue -- comes to rest against the hearth. The rest are lost in shadow.

\vspace{1em}
\noindent The Gray Wanderer picks up the slate. The unfinished game is a puzzle. \emph{"He was playing for time. He always played for time."}

\vspace{1em}
\noindent Outside, the Belworth churns. Riders are crossing. The road does not forgive a debt. It only reschedules.

\vspace{2em}
\noindent\textit{--- To be continued in "Blood on the Belworth."}

\clearpage
\chapter*{Appendix: The Canray Board (Optional Prop)}
The GM may use Braedon's unfinished Canray game as a framing device. The position of the pieces should reflect the political situation:

- **Blue (the throne):** Isolated, with few supporting pieces -- Veralyn's precarious position.
- **Oranges (allies):** Two forks -- one aligned with the Blue, one distant -- representing the divided loyalties of the Viterran and Vhasian courts.
- **Reds (the Fenwood lands):** A defensive wall along the home rank, but with a single gap where the Belworth runs.
- **Green (a runner):** Locked behind enemy lines -- the message Braedon wants delivered.

The solution to the game (if Braedon had lived) would have been to sacrifice the Blue to free the Green -- a metaphor for dying so that the heir could live. The GM may reveal this if the players ask.

\clearpage
\chapter*{Colophon}
\lettrine{T}{his} book was written in the solar at Tarlington, on paper that had once been a grain manifest. The ink is lamp‑black cut with the ash of a burned confession. The binding is horsehair, tied with nine knots.

The ninth knot is not cut.

The old duke is dead. The road is still south. The river rises.

\vspace{1em}
\noindent\textit{--- The Gray Wanderer, at the edge of the Belworth, dawn of the first day of the Crisis.}

\end{document}
